% Part: history
% Chapter: biographies 
% Section: alfred-tarski

\documentclass[../../../include/open-logic-section]{subfiles}

\begin{document}

\olfileid[pt]{his}{bio}{tar}

\olsection{Alfred Tarski}

\olphoto{tarski-alfred}{Alfred Tarski}

Alfred Tarski nasceu a 14 de janeiro de 1901 em Varsóvia, na Polónia
(então parte do Império Russo). Descrito como ``napoleónico'', Tarski
era ruidoso, falador e intenso.  A sua energia refletia-se muitas
vezes nas aulas---chegou a incendiar um cesto do lixo ao deitar fora
um cigarro durante uma aula, e foi proibido de lecionar novamente naquele
edifício.

Tarski tinha sede de conhecimento desde cedo. Mais tarde diria aos
estudantes que estudara lógica porque era a única disciplina em que
tivera um~B, mas os registros do liceu mostram que tinha notas
máximas em tudo---até em lógica. Estudou na Universidade de Varsóvia
de 1918 a 1924. Tarski pretendia inicialmente estudar biologia, mas
interessou-se por matemática, filosofia e lógica, pois a universidade
era o centro da Escola de Lógica e Filosofia de Varsóvia. Tarski
doutorou-se em 1924 sob orientação de Stanis\l{}aw Le\'{s}niewski.

Antes de emigrar para os Estados Unidos em 1939, Tarski completou
parte do seu trabalho mais importante enquanto professor de ensino
secundário em Varsóvia. O trabalho sobre consequência lógica e
verdade lógica foi escrito nesse período. Em 1939, Tarski encontrava-se
nos Estados Unidos numa digressão de conferências. Durante a visita,
a Alemanha invadiu a Polónia, e, por causa da ascendência judaica,
Tarski não pôde regressar.  A mulher e os filhos permaneceram na
Polónia até ao fim da guerra, mas depois também puderam emigrar para
os Estados Unidos.  Tarski lecionou em Harvard, no College of the City
of New York, no Institute for Advanced Study em Princeton, e
finalmente na Universidade da Califórnia, em Berkeley. Ali fundou o
programa multidisciplinar em Lógica e Metodologia da Ciência.  Tarski
morreu a 26 de outubro de 1983, aos 82 anos.

\begin{reading} 
Para mais sobre a vida de Tarski, ver a biografia \emph{Alfred Tarski:
  Life and Logic} \citep{Feferman2004}. As obras fundadoras de Tarski
sobre consequência lógica e verdade estão disponíveis em inglês em
\citep{Tarski1983}. Todas as obras originais de Tarski foram reunidas
numa série de quatro volumes, \citep{Tarski1981}.
\end{reading}

\end{document}
